O EDBank é uma aplicação web que simula operações bancárias para uso
em atividades de matemática financeira no Ensino Médio. A aplicação
organiza as contas em agências e permite aos participantes movimentar
valores fictícios, acompanhar registros no extrato e consultar as
consequências financeiras de suas operações.

A concepção do sistema reúne dois tipos de interação. Na área da
conta, o participante realiza operações como pagamentos,
transferências e compras. Na área de gerência, é possível configurar
a agência, acompanhar suas contas e executar ciclos de processamento.
Entre as operações realizadas em um ciclo está o crédito de salário
fictício, associado à profissão atribuída no momento da abertura da
conta.

\section{Público-alvo e Contexto de Uso}
\label{cap-cencepcao-publico}

O EDBank destina-se a estudantes do Ensino Médio em atividades de
matemática financeira. Em uma atividade orientada pelo professor,
o estudante pode prever o saldo após uma sequência de operações,
conferir o resultado no extrato e explicar como os pagamentos, os
créditos e as transferências afetaram o valor disponível.

A aplicação oferece uma área de conta para os participantes e uma
área de gerência para administrar agências. Também há previsão de
subgerência por convite e uma área administrativa do sistema.
Em um possível cenário de uso em sala de aula, o professor pode
utilizar a gerência de uma agência para organizar a atividade e
acompanhar as contas vinculadas a ela. A agência é a unidade de
organização adotada pela aplicação.

\section{Escopo da aplicação}
\label{cap-concepcao-escopo}

O EDBank compreende a abertura de contas fictícias em agências,
a consulta de saldo e extrato, as transferências entre contas,
o pagamento de cobranças, o uso de cartões, os investimentos
simulados e as compras em um shopping interno. Na área de gerência,
o sistema reúne funções relacionadas à configuração da agência,
à administração das contas e ao processamento de ciclos.

As movimentações utilizam valores fictícios e são processadas no
ambiente da aplicação, sem liquidação financeira real. A
autenticação com conta Google é utilizada para identificar o
participante no acesso ao sistema.

O módulo de empréstimos tem caráter demonstrativo. O cálculo está
previsto na interface, e a ação de solicitação exibe um protocolo
no navegador. Na versão analisada, essa ação não gera contrato,
registro de parcelas ou crédito em conta. A interface também inclui
recursos relacionados ao Pix, mas o fluxo de envio está incompleto
na versão analisada.


\section{Requisitos funcionais}
\label{cap-concepcao-requisitos-funcionais}

Os requisitos funcionais a seguir descrevem as ações previstas
para os participantes e para a gerência, a partir dos fluxos
identificados no código da aplicação.

\renewcommand{\arraystretch}{1.4}

\begin{xltabular}{\textwidth}{|>{\centering\arraybackslash}m{1.6cm}|>{\centering\arraybackslash}m{3.6cm}|X|}
\caption{Requisitos Funcionais do Sistema}\label{tab:requisitos-funcionais}\\
\hline
\textbf{Código} & \textbf{Item} & \textbf{Descrição} \\
\hline

RF01 & Autenticação & Permitir o acesso do participante por meio de uma conta Google. \\
\hline
RF02 & Abertura de conta & Permitir que o participante abra uma conta fictícia em uma agência, mediante código de acesso quando exigido, e associe a conta a uma profissão sorteada. \\
\hline
RF03 & Consulta de saldo e extrato & Apresentar o saldo da conta e o histórico de movimentações simuladas. \\
\hline
RF04 & Transferência entre contas & Permitir que o participante identifique uma conta de destino e transfira para ela um valor fictício. \\
\hline
RF05 & Pagamento de cobranças & Permitir a consulta de cobranças pendentes e seu pagamento com saldo em conta ou cartão, conforme as condições da operação. \\
\hline
RF06 & Cartões e faturas & Permitir a solicitação de cartão e a consulta de limite e fatura, além do pagamento da fatura no ambiente simulado. \\
\hline
RF07 & Investimentos & Permitir a aplicação e a retirada de valores fictícios, com consulta aos rendimentos apresentados pela aplicação. \\
\hline
RF08 & Compras no shopping & Permitir a consulta de produtos e a realização de compras fictícias com saldo ou cartão. \\
\hline
RF09 & Simulação de empréstimo & Permitir a estimativa de prestações a partir do valor solicitado e do prazo. O módulo possui interface e rotina de cálculo no navegador, cujo funcionamento precisa ser verificado em execução. \\
\hline
RF10 & Gerência de agência & Permitir a configuração da agência, a administração das contas vinculadas e o processamento de ciclos, incluindo o crédito de salário fictício. \\
\hline
\end{xltabular}

\section{Requisitos não funcionais}
\label{cap-concepcao-requisitos-nao-funcionais}

\section{Regras de negócio}
\label{cap-concepcao-regras}

\subsection{Contas, saldo e extrato}
\label{cap-concepcao-contas}

\subsection{Pagamentos e transferências simuladas}
\label{cap-concepcao-pagamentos}

\subsection{Simulação de crédito}
\label{cap-concepcao-credito}